%==============================================================================
\chapter{المرجعية المهنية وحُسن ممارسات تربية النحل}
%==============================================================================
يحدّد هذا الفصل المعارف المهنية التي تؤسّس قواعد التسيير والعتبات القابلة للضبط
ومؤشّرات النظام. يشكّل الأساس الوظيفي الذي تستند إليه لوحات القيادة والتنبيهات.

\section{عوامل إنتاج العسل}
يعتمد إنتاج الطائفة مباشرةً على الطقس وتوفّر غبار الطلع وتدفّق الرحيق. عندما يكون
تدفّق الرحيق مرتفعًا، تضع الملكة بيضًا أكثر، ويبلغ العدد ذروته، ويصل المردود
أقصاه. في فترة نقص الغذاء، يتراجع وضع البيض والعدد. لذلك يجب أن يربط النظام
الإنتاج المُقاس بالفترات والمواسم، ما يبرّر تصفية لوحات القيادة حسب الشهر
والأسبوع والموسم والسنة.

\section{موقع الموقع وإعداده}
يحدّد اختيار موقع الموقع الإنتاجية. يمكن للقيم المرجعية التالية أن تغذّي الحقول
وحدود ضبط الإدخال في وحدة إدارة المواقع.
\begin{xltabular}{\textwidth}{>{\bfseries}p{5.5cm} X}
\toprule
\textbf{المعيار} & \textbf{القيمة المرجعية} \\
\midrule
\endhead
القرب من الموارد & مصادر الرحيق وغبار الطلع على بُعد 1 كم تقريبًا. \\
الوصول إلى الماء & ماء صالح للشرب متاح، عدة لترات مُستهلَكة يوميًا لكل طائفة. \\
البُعد عن المساكن & 100 م في المنطقة الغابية، 200 م في المنطقة الشجيرية، 300 م في المنطقة المفتوحة. \\
البُعد عن المزروعات المعالَجة & 3 إلى 4 كم على الأقل عن المزروعات التقليدية التي تستعمل المبيدات. \\
الارتفاع والميلان & خلايا مرفوعة نحو 1.5 م عن الأرض ومائلة قليلًا لتصريف الرطوبة. \\
الحماية & مأوى من الشمس الحارقة والرياح والفيضانات. \\
\bottomrule
\end{xltabular}

\begin{figure}[htbp]
\centering
\includegraphics[width=0.48\linewidth]{hive-003}
\caption{مثال لخلية مركّبة على موقع: رفع على أرجل ولوحة طيران.}
\label{fig:ruche-installee}
\end{figure}

\section{أنواع الخلايا والمردود المرجعي}
\begin{xltabular}{\textwidth}{>{\bfseries}p{5.5cm} X}
\toprule
\textbf{نوع الخلية} & \textbf{الخصائص والمردود} \\
\midrule
\endhead
خلية تقليدية بأقراص ثابتة & 6 إلى 10 كغ من العسل في الأقراص لكل موسم، غير مُوصى بها في تربية النحل البيولوجية. \\
خلية بقضبان علوية & عرض معياري من 32 إلى 35 سم، تفتيش مستقلّ للأقراص. \\
خلية بإطارات (Langstroth، شرق-إفريقية) & أقراص عسل قابلة لإعادة الاستعمال عدة مواسم بعد الحصاد. \\
\bottomrule
\end{xltabular}
تُستعمل هذه المقادير التقريبية أساسًا لضبط عتبات الإنتاج وتقدير العائد على
الاستثمار لكل خلية.

\section{تصنيف نماذج الخلايا}
يُميَّز بين عائلتين كبيرتين من الخلايا، ويجب أن تبقى نمذجتهما منفتحة لتغطية
الاستعمالات المحلية والدولية.
\begin{xltabular}{\textwidth}{>{\bfseries}p{4.5cm} X}
\toprule
\textbf{العائلة} & \textbf{النماذج التمثيلية} \\
\midrule
\endhead
خلايا بإطارات (عمودية) & Dadant، Langstroth، Voirnot، Layens، Aumônière، Nationale (GB)، WBC، Tonelli. \\
خلايا تقليدية بلا إطارات & Warré، الكينية TBH، التنزانية TBH، جذع، قشّ. \\
نماذج إقليمية أخرى & Lucitana، Oksman، Smith، CDB، Duvauchelle، Jarry، Congres، Zanders، النماذج السلوفينية. \\
\bottomrule
\end{xltabular}
في فرنسا، أكثر النماذج شيوعًا هي خلايا Dadant وLangstroth وVoirnot. تتشارك جميع
الخلايا بإطارات تركيبًا مشتركًا، بأبعاد دقيقة إلى المليمتر، ما يعزّز نموذج
المعطيات المعتمَد.
\begin{encadre}
\textbf{إرساء النموذج:} النموذج المرجعي للمشروع هو خلية Dadant، التي يستوعب
جسمها حتى 10 إطارات. تؤسّس هذه القيمة السعة القصوى البالغة 10 إطارات لكل جسم أو
عاسلة المعتمَدة في قواعد التسيير. يبقى النظام مع ذلك قابلًا للضبط لدعم نماذج أخرى.
\end{encadre}

\subsection{التركيب المشترك لخلية بإطارات}
\begin{xltabular}{\textwidth}{>{\bfseries}p{4cm} X}
\toprule
\textbf{العنصر} & \textbf{الدور} \\
\midrule
\endhead
الأرضية / لوحة الطيران & قاعدة الخلية ومنطقة هبوط النحل. \\
الجسم (القاعدة) & مقصورة الحضنة الرئيسية، تحتوي على الإطارات القاعدية. \\
العاسلة & طابق تمديد مُخصّص لحصاد العسل. \\
الإطار & حامل متحرّك لأقراص الشمع، مركّب على موضع محدّد. \\
حاجز الملكة & يفصل الجسم عن العاسلات ويحصر الملكة في الجسم. \\
غطاء الإطارات والسقف & يُغلقان الخلية ويحميانها من التقلّبات الجوية. \\
المُغذّي & يتيح تزويد الطائفة بالغذاء. \\
\bottomrule
\end{xltabular}

\begin{figure}[htbp]
\centering
\includegraphics[width=0.85\linewidth]{hive-001}
\caption{مقطع لخلية بإطارات ودور كل عنصر (السقف، المُغذّي، غطاء الإطارات، العاسلة، حاجز الملكة، الجسم، الأرضية).}
\label{fig:coupe-ruche}
\end{figure}

\section{بروتوكول تفتيش الطوائف}
تُجرى عمليات التفتيش يُفضَّل في الأيام المشمسة وتتناول نقاط مراقبة دقيقة تُنظّم
تقرير الزيارة.
\begin{itemize}
  \item تطوّر الحضنة (البيض، اليرقات، الخلايا المُغطّاة)
  \item ملء الأعين بالعسل وغبار الطلع
  \item وجود طفيليات أو أمراض
  \item جمع الرحيق وغبار الطلع والعكبر
\end{itemize}
\begin{encadre}
\textbf{قيد تشغيلي:} يجب ألا تتجاوز الزيارة 45 دقيقة. بعدها، يزداد اضطراب الطائفة
وينتقل إلى الخلايا المجاورة. يمكن مراقبة حقل \emph{المدة} في تقرير الزيارة
مقارنةً بهذا الحدّ المرجعي.
\end{encadre}

\section{التطريد وتقسيم الطوائف}
يُنذر التطريد الوشيك بوجود خلايا ملكية على حواف الأقراص. قبل أيام من خروج ملكة
جديدة، تغادر الملكة القديمة الطائفة مع نصف الشغّالات. لذلك يجب أن يتيح النظام
استباق التقسيم. تُقرَّر ثلاث طرق مرجعية كأسباب تدخّل.
\begin{xltabular}{\textwidth}{>{\bfseries}p{4.5cm} X}
\toprule
\textbf{الطريقة} & \textbf{المبدأ} \\
\midrule
\endhead
تكوين نواة & نزع 1 إلى 2 قرص حضنة مع النحل إلى خلية جديدة. \\
نواة ملكة & نقل الملكة مع قرص خالٍ من الخلايا الملكية. \\
تطريد اصطناعي & نقل الخلية الأمّ ووضع خلية جديدة في الموضع القديم مع 2 إلى 3 أقراص حضنة يافعة. \\
\bottomrule
\end{xltabular}

\section{مؤشّرات الرفاه وأسباب الانخفاض}
تُظهر الطائفة السليمة عشّ حضنة متطوّرًا جيدًا في مختلف الأطوار، وأعينًا مملوءة
بالموارد، ونشاط رعي مرئيًا، وغياب الطفيليات والأمراض. تغذّي هذه المعايير التقييم
النوعي للحالة الصحية ومستوى الإنتاجية المُقيَّم من 1 إلى 3 في تقرير الزيارة.

ينتج هروب الخلية أو هجرها عن اضطرابات مفرطة أو ظروف غير ملائمة (نقص الغذاء أو
الماء، حرارة قصوى). تركُ احتياطي من العسل للطائفة عند الحصاد يقلّل هذا الخطر.
تبرّر هذه الظواهر تنبيهات الانخفاض المفاجئ في العدد أو الإنتاج والعتبات القابلة
للضبط المرتبطة بها.

\section{الأثر على العتبات القابلة للضبط}
\begin{xltabular}{\textwidth}{>{\bfseries}p{4.5cm} X}
\toprule
\textbf{معلمة النظام} & \textbf{الإرساء المهني} \\
\midrule
\endhead
عتبة جمع العسل & المردود المرجعي حسب نوع الخلية والموسم. \\
عتبة تنبيه الانخفاض & انخفاض غير طبيعي مرتبط بالهروب أو المرض أو نقص الغذاء. \\
المدة القصوى للزيارة & حدّ 45 دقيقة لكل تفتيش. \\
مسافات أمان الموقع & 100 أو 200 أو 300 م حسب المنطقة، 3 إلى 4 كم عن المزروعات المعالَجة. \\
فترات التحليل & ربط الإنتاج بالموسم (تدفّق الرحيق). \\
\bottomrule
\end{xltabular}
